% =====================================================================
%  On the Necessary Substructures of Finite Contact Graphs
%  A conditional-existence account of why a finite world is structured.
% =====================================================================
\documentclass[11pt,a4paper]{article}

\usepackage[utf8]{inputenc}
\usepackage[T1]{fontenc}
\usepackage{lmodern}
\usepackage[a4paper,margin=2.5cm]{geometry}
\usepackage{amsmath,amssymb,amsthm,mathtools}
\usepackage{bm}
\usepackage{mathrsfs}
\usepackage{booktabs}
\usepackage{array}
\usepackage{enumitem}
\usepackage{microtype}
\usepackage{graphicx}
\usepackage[round,sort&compress]{natbib}
\usepackage{tikz}
\usetikzlibrary{arrows.meta,positioning,calc,decorations.pathmorphing,fit}
\usepackage[colorlinks=true,linkcolor=blue!45!black,
            citecolor=blue!45!black,urlcolor=blue!45!black]{hyperref}
\usepackage{cleveref}

% ---- Theorem environments ----
\newtheoremstyle{thmstyle}{\topsep}{\topsep}{\itshape}{}{\bfseries}{.}{.5em}{}
\theoremstyle{thmstyle}
\newtheorem{theorem}{Theorem}[section]
\newtheorem{lemma}[theorem]{Lemma}
\newtheorem{proposition}[theorem]{Proposition}
\newtheorem{corollary}[theorem]{Corollary}
\theoremstyle{definition}
\newtheorem{definition}[theorem]{Definition}
\newtheorem{axiom}{Axiom}
\newtheorem{principle}[theorem]{Principle}
\newtheorem{construction}[theorem]{Construction}
\newtheorem{example}[theorem]{Example}
\theoremstyle{remark}
\newtheorem{remark}[theorem]{Remark}
\newtheorem{notation}[theorem]{Notation}

\crefname{axiom}{Premise}{Premises}
\Crefname{axiom}{Premise}{Premises}
\crefname{principle}{Principle}{Principles}
\Crefname{principle}{Principle}{Principles}

% ---- Shorthand ----
\newcommand{\R}{\mathbb{R}}
\newcommand{\N}{\mathbb{N}}
\newcommand{\Z}{\mathbb{Z}}
\newcommand{\Gph}{\Gamma}                 % the contact graph
\newcommand{\V}{V}                        % vertex set
\newcommand{\E}{E}                        % edge set
\newcommand{\wt}{w}                       % edge weight
\newcommand{\floorc}{\beta}               % the floor
\newcommand{\Nbr}{N}                      % neighbourhood
\newcommand{\co}{\complement}             % complement
\newcommand{\cut}{\partial}               % edge boundary / cut
\newcommand{\Walk}{W}                     % a walk
\newcommand{\gate}{g}                     % gating function
\newcommand{\Agent}{A}                    % agent subgraph
\newcommand{\Res}{\rho}                   % residual invariant
\newcommand{\Quot}{Q}                     % quotient graph
\newcommand{\rec}{M}                      % committed record
\newcommand{\Ver}{\mathsf{V}}             % verifier
\newcommand{\abs}[1]{\lvert#1\rvert}
\newcommand{\norm}[1]{\lVert#1\rVert}
\newcommand{\restr}{\!\restriction\!}

\title{\textbf{On the Necessary Substructures of Finite Contact Graphs:\\[2pt]
A Conditional-Existence Account of Why a Finite World Is Structured}}

\author{%
Kundai Farai Sachikonye\\
\small Technical University of Munich\\
\small AIMe Registry for Artificial Intelligence\\
\small \texttt{kundai.sachikonye@wzw.tum.de}}

\date{\today}

\begin{document}
\maketitle

% =====================================================================
\begin{abstract}
\noindent
We ask a structural question---\emph{why is a world structured at all?}---and
answer it with a single hypothesis about graphs. A \emph{contact graph} is a
finite weighted graph whose vertices are the parts that have been individuated
within a whole and whose edges carry the cost of separating one part from
another. We take as our only premise that such a graph is \emph{finite} and
that its edge weights are bounded below by a positive constant; we call this the
Finiteness Premise. From it alone we derive, as a chain of existence theorems, a
ladder of substructures that any finite contact graph must contain. We prove:
\textbf{(T0)} a strictly positive \emph{floor} on every separating boundary---no
two parts are adjacent at zero cost; \textbf{(T1)} that each vertex is fixed only
by its complement-neighbourhood, so that individuation must proceed by negation
and admits no selector without regress; \textbf{(T2)} that the
\emph{cut-residual}---the inter-part boundary content---is bounded below and is
invariant under every relabelling (isomorphism) of the graph, while no single
vertex realises it, so the residual is a region and never a point; \textbf{(T3)}
that this residual is the conserved quantity under all re-encodings of the
graph's vertices, the unique invariant from which the relabellable surface
derives; \textbf{(T4)} that every connected subgraph carries its own such
invariant under its automorphisms; \textbf{(T5)} that selecting a walk on the
graph requires an edge-gating function, without which the next step is
undetermined; \textbf{(T6)} that the committed-edge count along any walk is
monotone, so revisiting a vertex is never returning to a state and resemblance
never entails identity; \textbf{(T7)} that a coherent family of subgraphs acting
as one admits a single gating function on their quotient graph, and only one;
and \textbf{(T8)} that, because the subgraph invariant is complete, contact
graphs are specifiable and hence instantiable---authored finite structures
exist---subject to three bounds we prove: their walks are not foreseeable from
their specification, an authored copy is a distinct structure, and authorship is
limited to the resolution of the floor. We close with a Master Theorem: every
substructure T0--T8 is forced for finite graphs and is \emph{not} forced for
infinite or zero-floor graphs. Structure is therefore neither imposed on a world
nor intrinsic to it independent of its parts; it is the necessary form a world
must take to be resolved by finitely many bounded parts. \emph{Structure is the
shadow of finitude.} Every claim is proved from the single premise; uncertain
material has been omitted, and the interpretations a reader may attach are
confined to clearly marked remarks on which no theorem depends.
\end{abstract}

\noindent\textbf{Keywords:} contact graph; finite weighted graph; edge-weight
floor; individuation by negation; cut invariant; graph isomorphism invariant;
gated walk; monotone walk; quotient graph; conditional existence; necessary
substructure; authored structure.

\tableofcontents

% =====================================================================
\section{Introduction}
\label{sec:intro}
% =====================================================================

\subsection{The question}

Why is a world structured at all? One might expect a world to be either a
featureless continuum, in which nothing is marked off from anything else, or an
arbitrary heap, in which everything is marked off from everything else at no
cost. Neither is what is observed: parts are distinguished, but their
distinction is bounded---there is a cost to telling one part from another, and
the parts compose into a definite architecture. We argue that this architecture
is not contingent. It is forced, in full, by a single fact: that the parts are
\emph{finite in number} and \emph{finitely separated}. The structure of a world
is the shadow its finiteness casts.

We make this precise in the language of graphs, because every relation of
``this part, separated from that part'' is natively an edge between two
vertices, and any architecture of parts is some graph. The object of the paper
is therefore a single graph---a \emph{contact graph}---whose vertices are the
individuated parts of a whole and whose weighted edges record the cost of the
separations between them. We pose one hypothesis about this graph and derive
everything else.

\subsection{Method: conditional existence}

The paper proves no empirical claim. It establishes a conditional: \emph{if} a
finite contact graph exists, \emph{then} a determinate ladder of substructures
must exist within it. Each rung is an existence theorem of the form ``under the
Finiteness Premise, the following is forced.'' We assume nothing about what the
vertices \emph{are}; we assume only that they are finite in number and
positively separated, and we read off what their graph cannot lack. The
deductions are combinatorial and order-theoretic throughout. Where a tempting
claim could not be proved at this standard it has been removed rather than
weakened.

\subsection{The single premise}

\begin{axiom}[Finiteness Premise]
\label{ax:finite}
A \emph{contact graph} is a finite weighted graph $\Gph=(\V,\E,\wt)$ with
$\abs{\V}<\infty$, in which every edge weight is bounded below by a positive
constant: there is $\floorc>0$ with $\wt(e)\ge\floorc$ for all $e\in\E$. We
assume throughout that a contact graph exists and that it is connected.
\end{axiom}

The premise is almost nothing: a finite graph whose edges cost something. Its
consequences are almost everything we mean by structure.

\subsection{The ladder of necessary substructures}

Under \Cref{ax:finite} we prove the following chain, each rung from the premise
together with those above it.

\begin{enumerate}[label=\textbf{(T\arabic*)},leftmargin=3em,start=0]
\item \emph{The floor exists} (\Cref{sec:floor}): every separating boundary has
weight $\ge\floorc>0$; no two parts are adjacent at zero cost.
\item \emph{Individuation is by negation} (\Cref{sec:negation}): a vertex is
fixed only as the complement of its non-neighbours; no selector fixes it without
regress.
\item \emph{Truth is a residual region} (\Cref{sec:residual}): the cut-residual
is bounded below and realised by no single vertex, hence a region, never a
point.
\item \emph{The residual is the conserved invariant} (\Cref{sec:residual}):
under every relabelling of vertices the residual is preserved; it is the unique
invariant from which the relabellable surface derives.
\item \emph{Each subgraph has a private invariant} (\Cref{sec:invariant}): every
connected subgraph carries its own invariant under its automorphisms.
\item \emph{Selection requires a gate} (\Cref{sec:gating}): a walk is
undetermined without an edge-gating function that biases the next step.
\item \emph{Histories do not return} (\Cref{sec:monotone}): the committed-edge
count is monotone; revisiting a vertex is not returning to a state.
\item \emph{Plurality forces a single quotient gate} (\Cref{sec:quotient}): a
coherent family of subgraphs acting as one admits exactly one gating function on
their quotient.
\item \emph{Authored structures exist, and are bounded} (\Cref{sec:construct}):
the subgraph invariant is complete, so contact graphs are specifiable and
instantiable; their walks are unforeseeable from the specification, an authored
copy is a distinct structure, and authorship reaches only to the floor.
\end{enumerate}

\Cref{sec:master} proves the Master Theorem: T0--T8 are forced for finite
graphs and not forced when finiteness or the positive floor is dropped.

\subsection{On the word ``contact''}

We use \emph{contact} in one operational sense: two parts are \emph{in contact}
when the graph carries an edge between them, and the weight of that edge is the
cost of the separation that individuates each from the other. Nothing in the
paper depends on any informal reading of the word; ``contact'' names an edge and
its weight, and the theorems about contact graphs are theorems about finite
weighted graphs in exactly this sense.

\subsection{Why finiteness is the source}

Finiteness is not a convenience of presentation but the origin of every
structure below. On an infinite graph a vertex may be separated from the rest by
edges of vanishing weight, the residual may be driven to zero, and a walk may be
specified to unbounded depth at fixed cost; none of the substructures is then
forced. It is the finitude of the vertex set, together with the positive floor
on edge weights, that forbids the zero-cost separation and compels the ladder.
We make the premise exactly strong enough to force the structures and no
stronger, and we mark, at each rung, where finiteness is used
(\Cref{sec:master}).

\subsection{Relation to existing work}

The paper is pure mathematics. It uses graph theory~\citep{bondymurty2008,%
diestel2017,west2001,bollobas1998}, algebraic and spectral graph
theory~\citep{godsilroyle2001,chung1997,fiedler1973,cheeger1970}, the theory of
walks and Markov chains on graphs~\citep{lovasz1993,levinperes2017}, order and
fixed-point theory~\citep{daveypriestley2002,birkhoff1967,tarski1955},
enumerative combinatorics and partitions~\citep{stanley2012,andrews1976},
self-reference and diagonal arguments~\citep{russell1903,cantor1891,lawvere1969,%
godel1931}, and information theory~\citep{shannon1948,cover2006}. Where the
development meets the sorites and the logic of vagueness we note the
connection~\citep{black1937,williamson1994} without taking any position in that
literature; our treatment is geometric, not a logic of partial truth. A
reasonable reader will recognise that the abstract premise is satisfied by a
range of concrete systems of finitely many bounded parts; we make no such
identification and prove everything of the abstract graph alone. Such
interpretive readings as we offer are confined to remarks labelled accordingly,
on which no theorem depends.

\subsection{Organisation}

\Cref{sec:setting} fixes the contact graph and its derived objects.
\Cref{sec:floor} proves T0. \Cref{sec:negation} proves T1. \Cref{sec:residual}
proves T2 and T3. \Cref{sec:invariant} proves T4. \Cref{sec:gating} proves T5.
\Cref{sec:monotone} proves T6. \Cref{sec:quotient} proves T7.
\Cref{sec:construct} proves T8. \Cref{sec:master} proves the Master Theorem and
concludes.
